\chapter{Basic Configuration --- pfSense Firewall}

\section{Firewall -- pfSense CE 2.6 Setup}

The firewall configuration in this project was implemented by following the step-by-step procedure shown in the video below, adapted to the EVE-NG lab environment.

\subsection{Reference Video}

\begin{itemize}
\tightlist
\item
  \url{https://youtu.be/sX22a_v2svQ?si=IzMju0KYcMQSE4sZ}
\end{itemize}

\begin{center}\rule{0.5\linewidth}{0.5pt}\end{center}

\subsection{Phase 1 -- Image preparation (EVE-NG SSH as root)}

pfSense CE 2.6 requires a disk image before the node can be started. Unlike VyOS, pfSense installs itself to disk during the first-boot wizard --- no manual \texttt{install\ image} step needed.

\begin{lstlisting}
cd /opt/unetlab/addons/qemu/pfsense-ce-2.6.0-release-amd64/

# List contents -- you should see the ISO
ls
# cdrom.iso

# Create the target disk (8G is enough for CE 2.6)
qemu-img create -f qcow2 virtioa.qcow2 8G

# Fix permissions
cd /opt/unetlab/addons/qemu/
/opt/unetlab/wrappers/unl_wrapper -a fixpermissions
\end{lstlisting}

\begin{lstlisting}
The disk preparation is identical to the VyOS flow.
See [Router (VyOS)](Router.md) for reference.
\end{lstlisting}

\begin{center}\rule{0.5\linewidth}{0.5pt}\end{center}

\subsection{Phase 2 -- First boot and setup wizard (VNC console)}

Start the node from the EVE-NG web UI and open the console (VNC).

pfSense will boot from the ISO and launch the installer automatically.

\subsubsection{2.1 Install to disk}

When prompted, accept the default options:

\begin{itemize}
\tightlist
\item
  \textbf{Install pfSense} $\rightarrow$ Continue
\item
  \textbf{Auto (ZFS)} or \textbf{Auto (UFS)} --- UFS is simpler for a lab
\item
  Select disk: \texttt{vtbd0} (the \texttt{virtioa.qcow2} disk)
\item
  Confirm and proceed --- pfSense will copy files and reboot
\end{itemize}

After reboot the node boots from disk. The ISO is no longer needed but can stay.

\subsubsection{2.2 Interface assignment}

On first boot, pfSense asks which physical interfaces to assign as WAN and LAN.

pfSense sees interfaces in the order EVE-NG/QEMU presents them:

{\footnotesize\begin{longtable}[]{@{}
  >{\raggedright\arraybackslash}p{0.22\linewidth}
  >{\raggedright\arraybackslash}p{0.22\linewidth}
  >{\raggedright\arraybackslash}p{0.45\linewidth}@{}}
\toprule\noalign{}
\textbf{pfSense NIC} & \textbf{EVE-NG port} & \textbf{Connected to} \\
\midrule\noalign{}
\endhead
\bottomrule\noalign{}
\endlastfoot
\texttt{vtnet0} (em0) & e0 & LAN network ($\rightarrow$ Router) \\
\texttt{vtnet1} (em1) & e1 & WAN network (pnet / internet) \\
\end{longtable}}

Always verify the assignment matches the lab cabling before continuing.

At the console menu:

\begin{lstlisting}
Should VLANs be set up now? $\rightarrow$ n
Enter the WAN interface name: vtnet1
Enter the LAN interface name: vtnet0
Do you want to proceed?      $\rightarrow$ y
\end{lstlisting}

pfSense will assign:

\begin{itemize}
\tightlist
\item
  \textbf{WAN} (\texttt{vtnet1}): DHCP by default --- gets IP from the upstream network
\item
  \textbf{LAN} (\texttt{vtnet0}): \texttt{192.168.1.1/24} by default --- \textbf{change this to match your lab}
\end{itemize}

\subsubsection{2.3 Set LAN IP from console}

From the pfSense console menu, select option \texttt{2)\ Set\ interface(s)\ IP\ address}:

\begin{lstlisting}
Select interface to configure: 2 (LAN)
Enter the new LAN IPv4 address:  172.16.x.x
Enter the new LAN IPv4 subnet:   30
For a WAN, enter the upstream gateway: (blank)
Do you want to enable DHCP on LAN? n
\end{lstlisting}

\begin{center}\rule{0.5\linewidth}{0.5pt}\end{center}

\subsection{Phase 3 -- Web GUI access}

pfSense web GUI is only reachable from the \textbf{LAN interface}.

{\footnotesize\begin{longtable}[]{@{}>{\raggedright\arraybackslash}p{0.45\linewidth}
  >{\raggedright\arraybackslash}p{0.45\linewidth}@{}}
\toprule\noalign{}
Parameter & Value \\
\midrule\noalign{}
\endhead
\bottomrule\noalign{}
\endlastfoot
URL & \texttt{https://172.16.x.x} \\
Default user & \texttt{admin} \\
Default password & \texttt{pfsense} \\
\end{longtable}}

\begin{lstlisting}
The web GUI and SSH are blocked on the WAN interface by default.
Connect from a host that has L2 reachability to the LAN segment
(e.g. via the Router or a management host on the same bridge).
\end{lstlisting}

Complete the setup wizard on first login:

\begin{enumerate}
\def\labelenumi{\arabic{enumi}.}
\tightlist
\item
  Set hostname and DNS
\item
  Set timezone
\item
  Confirm WAN settings
\item
  Confirm LAN IP
\item
  Change admin password
\item
  Reload --- pfSense is ready
\end{enumerate}

\begin{center}\rule{0.5\linewidth}{0.5pt}\end{center}

\subsection{Phase 4 -- Enable SSH management}

Go to \textbf{System $\rightarrow$ Advanced $\rightarrow$ Admin Access}:

\begin{itemize}
\tightlist
\item
  Enable \textbf{Secure Shell}
\item
  Set \textbf{SSHd Key Only} to \emph{Password or Public Key} for lab convenience
\item
  Click \textbf{Save}
\end{itemize}

SSH is then available from the LAN segment:

\begin{lstlisting}
ssh admin@172.16.x.x
\end{lstlisting}

\begin{center}\rule{0.5\linewidth}{0.5pt}\end{center}

\subsection{Phase 5 -- Persistence check}

pfSense writes its config to \texttt{/cf/conf/config.xml} on disk. After any change via web GUI or console, it saves automatically.

To verify the disk is being used:

\begin{lstlisting}
# From pfSense console $\rightarrow$ option 8 (Shell)
df -h
# /dev/vtbd0p3 should be mounted at /
\end{lstlisting}

\begin{quote}
Use \textbf{Halt System} from the pfSense menu or \texttt{shutdown -h now} from the shell. A hard stop from EVE-NG GUI may corrupt the ZFS/UFS filesystem.
\end{quote}

\begin{center}\rule{0.5\linewidth}{0.5pt}\end{center}

\subsection{Key notes from lab experience}

\begin{itemize}
\tightlist
\item
  pfSense console is \textbf{VNC only} in EVE-NG (no serial/telnet access from host).
\item
  SSH and web GUI are \textbf{disabled on WAN} by default --- always connect via LAN.
\item
  The tap interface connecting pfSense WAN to the \texttt{pnet} bridge may need to be re-added manually after each EVE-NG host reboot (see \href{../../labs/lab1/firewall.md}{LAB1 Firewall}).
\end{itemize}
